% V22: French physical pp.85 (Grosso modo)--88 (first paragraph);
% English physical pp.55 (Grosso modo)--57 (before differential properties).
粗略地说，譬如假设$f$是浸入，那么一般锥投影的性质，会因假设$\dim X\geq m$还是$\dim X\leq m$，甚至$\dim X<m$，而有很大不同。\cmt{V22-FIB001}{法85末至86首为“voire $\dim X<m$”；英56将voire误写成“see”，旧译因而写“另见”。这里恢复“甚至”。原稿保留两个在等号处相交的范围，本句是概括性的引导，不是互斥分类。}
以下考虑与$\operatorname{Grass}^{m+1}(P)$的泛点$\eta$对应的
$C_\eta\subset P_{k(\eta)}$，不再把所得结果逐一解释成关于“几乎所有点……”的说法。

首先，我们已在5.3中指出（第12节补述了一般情形），$C_\eta$与$X_{k(\eta)}$“正则”相交。更确切、也更一般地，对于$X$上的每个拟凝聚层$F$，在$X_{k(\eta)}$上相应的秩为$m+1$的局部自由模的截面$\varphi_\eta^{(m+1)}$，在Koszul意义下是$F_{k(\eta)}$-正则的；其零点子概形是$f_{k(\eta)}^{-1}(C_\eta)$。\cmt{V22-FIB002}{法86的截面在$X_{k(\eta)}$上，却将其零点概形简写成$C_\eta$；正文明确写出逆像。法稿局部自由模的秩清楚写$m+1$，英56误写$n+1$；法稿上下标为$\varphi_\eta^{(m+1)}$，英稿又把$\eta$挪到零点概形后的文字中，均据法稿恢复。原稿未在此重述正则性的约定。本版对任意拟凝聚模核验的是Koszul复形无正次同调；若$F$凝聚，则在其支撑与零点的交上等价于局部坐标构成普通正则列。零点外不强令单位生成的商模非零，也不把非有限模的Koszul正则性无条件改称任意坐标顺序下的普通正则列。泛系数论证及有限型情形的爆破识别另见核查报告V22-FIB-E02。}
由10.2，例如可知，态射
\[
\widetilde X(C_\eta)\longrightarrow X_{k(\eta)}
\]
将$\widetilde X(C_\eta)$等同于沿
\[
f_{k(\eta)}^{-1}(C_\eta)
 =X_{k(\eta)}\times_{P_{k(\eta)}}C_\eta
\]
爆破$X_{k(\eta)}$所得到的预概形。
在$\dim f(X)\leq m$时，甚至有$f_{k(\eta)}^{-1}(C_\eta)=\varnothing$，因而
\[
\widetilde X(C_\eta)\xrightarrow{\sim}X_{k(\eta)}
\]
是同构（事实上，此时受限锥投影本来就定义在整个$X_{k(\eta)}$上）。\cmt{V22-FIB003}{法86的$\dim f(X)\leq m$可辨。这里以及下文的像维数，均指像的闭包的维数。泛中心避开维数$\leq m$的闭集；必须先换底到$k(\eta)$，不能把这一结论写成中心已在$k$上定义。法稿在关联空间的波浪号上有所省略，本版统一为$\widetilde X(C_\eta)$。}
接着便要问
\[
p_{C_\eta}:\widetilde X(C_\eta)\longrightarrow Q(C_\eta)
\]
的纤维维数及这个态射的平坦性。我们得到：
\cmt{V22-FIB004}{法86写“et de la platitude de ce morphisme”，是并列的研究问题；英56误接为“and we find the flatness”。正文恢复问句的范围，不预先宣称无条件平坦。以下维数论证置于本篇的域上有限型背景；一般无限维或非诺特对象不在这项订正命题的范围内。}

\egas{14.1}{命题〔条件订正〕}
假设$X$是域$k$上的有限型预概形。记$X_i$为它的不可约分支，$x_i$为$X_i$的泛点。
要求每个$X_i\to P$的所有非空几何纤维均为纯$d$维，且$d$与$i$无关。
例如，当$f:X\to P$拟有限时，这个条件成立，且$d=0$。\cmt{V22-FIB005}{法86原先仅假设$X$不可约，继而推广为各$X_i$在$f(x_i)$上的\emph{泛纤维}维数共同为$d$；仅此不足以保证(a)。取$X=\operatorname{Bl}_p(\mathbb P^3_k)$，$f$为爆破态射，$m=2$；$X$不可约且$d=0$，但通过泛中心点$C_\eta$及$p$的直线所对应的投影纤维含有例外$\mathbb P^2$，维数为2而非$\dim X-m=1$。正文明确加强到所有非空几何纤维纯$d$维，不冒充原稿已写出的前提。此充分版本及反例详见核查报告V22-FIB-E03、E04。}
则有：

\noindent(a) 若$\dim f(X)>m$，则
\[
p_{C_\eta}:\widetilde X(C_\eta)\longrightarrow Q(C_\eta)
\]
的每个纤维均非空，且维数等于$\dim X-m$。

\noindent(b) 若$\dim f(X)\leq m$，则$p_{C_\eta}$的每个非空纤维维数均为$d$。
特别地，若$f:X\to P$有限（分别地，拟有限），则$p_{C_\eta}$有限（分别地，拟有限）。\cmt{V22-FIB006}{法86左缘手写恢复了$\dim f(X)\leq m$；英56误作$\dim x\leq m$。法稿还补有“$X_i$在$P$上的非空纤维都为维数$d$”，并在$p_C$的纤维下面插入“non vides”。故(b)的“非空”不是本译者猜补。末尾有限／拟有限条件有删改；可辨的现行意思为$f$有限／拟有限时的相应结论。不能从任意$d$维纤维直接推出拟有限；在这两个特殊情形中$d=0$。被划去的“于是$p_C$拟有限”及“$f$在$k$上固有”等旧层另存源页与报告，不把删字接回现行句。}

\eproof
在(a)的情形，对$Q(C_\eta)$的每个点$\xi$，我们已有
\[
\dim X^{(m)}_\xi\geq\dim X-m
\]
（我希望如此！）。\cmt{V22-FIB007}{法86末至87首原称对$\operatorname{Grass}^m(P)$的\emph{所有}点均有此下界，再说特别适用于$Q(C_\eta)$。对非固有$X$或只有较小分支遇到截面的情况，前一无条件说法不能照用；正文只保留本命题所需且在订正条件下成立的$Q(C_\eta)$范围。最大维分支确实遇到每个相关截面，须用泛中心与像闭包边界的维数界；核查报告E04给出这一步。法稿的“je l'espère !”属于作者核对语，照存。}
对于反向不等式，置身于域$k'=k(\xi)$之上。
$C_{\eta,k'}\subset L_\xi$是$L_\xi$中的一个超平面。如果
\[
\dim X^{(m)}_\xi
 =\dim\bigl(X_{k'}\times_{P_{k'}}L_\xi\bigr)
 \geq\dim X-m+1,
\]
则中心截面的维数至少为$\dim X-m$。这里，基变换$k(\eta)\to k'$把中心截面变为
\[
\bigl(X_{k'}\times_{P_{k'}}L_\xi\bigr)
       \times_{L_\xi} C_{\eta,k'}.
\]
另一方面，由第2节（在第10节中复查过），却有
\[
\dim X^{(m+1)}_\eta=\dim X-m-1,
\]
这便矛盾。\cmt{V22-FIB008}{英56漏去了完整条件式$\dim X^{(m)}_\xi\geq\dim X-m+1$，又破坏了纤维积；法87两式均可辨，已据实恢复。这里的中心截面是$X^{(m+1)}_\eta=X_{k(\eta)}\times_{P_{k(\eta)}}C_\eta$，不是$X^{(m)}_\xi$。法稿在末式后划去了“至少假设$\dim f(X)>m$”；正文已有(a)的这一假设。非固有情形不能单凭“超平面”断言任意正维局部闭集必相交，故这条反证链须配合E04中的像闭包及边界论证；不能从恢复原排式就推定原弱前提已获证明。}

(b)用类似的方法处理：如果
\[
\dim\bigl(X_{k'}\times_{P_{k'}}L_\xi\bigr)\geq d+1,
\]
也就是说，
\[
\dim\bigl(f_{k'}(X_{k'})\cap L_\xi\bigr)\geq1,
\]
那么用与上面相同的论证便会得到
$X^{(m+1)}_\eta\ne\varnothing$，这与14.1之前指出的事实矛盾。\cmt{V22-FIB009}{法87明确写的是像与$L_\xi$的\emph{交集}；英56误排成“$f_{k'}(X_{k'})\to L$”，不是另一相对维数。正文依可辨法稿恢复。严格论证可在像的射影闭包中进行：它与中心不交，而与$L_\xi$的正维闭交必遇到$L_\xi$内的超平面$C_{\eta,k'}$，不可能。因此该闭交为有限集，再以$f$的纤维维数条件计算；详见E04。法稿本段原起句“$\dim f(X)=m$的情形也可由(b)处理”已被作者划去，现行起首改为“(b)的情形”。}

\egas{14.2}{推论}
假设$\dim X=m$，且$f:X\to P$有限（分别地，拟有限）。则态射
\[
p_{C_\eta}:X_{k(\eta)}\longrightarrow Q(C_\eta)
\]
有限满射（分别地，拟有限且优势）。

事实上，这个态射是拟有限的；又因为
\[
\dim X_{k(\eta)}=\dim Q(C_\eta),
\]
所以它是优势的。若$f$有限，则$p_{C_\eta}$还是有限的，因而固有；既然优势，便为满射。\cmt{V22-FIB010}{法87已将证明首句的“fini”手改为“quasi-fini”；英57此处保留了改后读法。命题中有限与拟有限的两种情况分别对应，不能把后一种加强成满射。完整的有限性理由为：泛中心避开像的射影闭包，闭包的受限投影是固有拟有限态射，故有限；再与$f$复合。优势性来自拟有限态射保持像维数及$\dim X=m$，不是任意两个等维对象之间的态射都优势。}

\egas{14.3}{推论〔补明分支条件〕}
在14.1(a)的订正条件下，再假设$X$的每个不可约分支$X_i$均满足$\dim f(X_i)\geq m$。
若$X$是Cohen--Macaulay的，则
\[
p_{C_\eta}:\widetilde X(C_\eta)\longrightarrow Q(C_\eta)
\]
是Cohen--Macaulay态射，因而特别是平坦的。\cmt{V22-FIB011}{法87原推论只写“在14.1(a)的条件下”。除14.1已补强的纤维条件外，可约情形还须防止低维像分支。例如$X=\mathbb P^2_k\amalg\operatorname{Spec}k\to\mathbb P^3_k$由一个平面及其外一点的嵌入给出，$m=1$。$X$是CM、$f$拟有限，且$\dim f(X)>1$，但孤立点分支到$Q(C_\eta)\simeq\mathbb P^1_{k(\eta)}$不平坦。正文所加各分支像维数条件排除此问题；在$f$拟有限时，它正是下文作者模版本已有的支撑分支维数条件。补证见E05。}

原稿给出的证明提示是：参看第21页上方、第5节之前的注，并将其更强的结果并入14.3。在紧随的CM凝聚模版本及其支撑维数条件下，可表述为：对$\xi\in Q(C_\eta)$，截面$\varphi_\xi^{(m)}$在其零点与模的支撑相交处是$F_{k(\xi)}$-正则的。\cmt{V22-FIB012}{法87手写为“à inclure dans14.3”，以及$\varphi_\xi^{(m)}$，英57将前者变成疑问、后者误排为$C_\xi^{(m)}$。法稿原句直接说该截面为$F$-正则，未在这句重述$F$的条件；不能据此将断言扩大到开头的任意拟凝聚模，例如低维支撑模便可能不满足所需正则列条件。正文明确限定到下述CM凝聚模版本，并标明此段是证明提示；零点外仍按Koszul复形可缩理解。法稿写模的底域$k(\eta)$，在$\xi$处作截面应再换底至$k(\xi)$，正文补明。所引“第21页上方、第5节之前”在当前法稿扫描中未能定位：物理第21页直接从第5节开始；英文编者也要求订正此处。故仍保留原交叉引用及其未解决状态，不能以此引用代替E05的独立补证。}

这个推论应当改写为模的版本；为简单起见，可以假设$f$拟有限：若$F$是$X$上的Cohen--Macaulay凝聚模，且对于$\operatorname{Supp}F$的每个不可约分支$Z$都有$\dim Z\geq m$，那么$\widetilde F(C_\eta)$相对于$Q(C_\eta)$是Cohen--Macaulay的，因而特别是平坦的。\cmt{V22-FIB013}{法87用“devrait être modulé”，接着立即改谈模$F$；英57笼统译为“must be modified”。正文恢复这一模版本指向，同时保留“应当”的编写语气。“凝聚”明列本版补证采用的有限模范围；不得将此CM模定理直接推广到任意非有限拟凝聚模。相对CM包括对基底平坦和各纤维上的CM性；这里不是只断言总空间上的模CM。}

注意，即使只要求$p_{C_\eta}$平坦，也不能把$X$上的CM假设仅仅换成一个维数假设。例如，假设$f$是浸入，而$X$不可约、维数为$m$，于是$p_{C_\eta}$拟有限。
由于$X_{k(\eta)}$与$Q(C_\eta)$不可约且维数相同，而后者正则，只有当$X_{k(\eta)}$是CM时，$p_{C_\eta}$才可能平坦。\cmt{V22-FIB014}{法88明确写“$X$ irr de dim $m$”，故英57的“$f$ is irreducible”是转写错误。此处$k(\eta)/k$是Grassmann泛点的纯超越扩张，保留不可约性；拟有限平坦态射在点上的零维纤维是CM，正则基底加上平坦CM纤维使源CM。报告E06另给出一个二维整非CM仿射概形的具体例子，以验证仅有维数条件确实不足。}

